\chapter{Geltungsbereich und metrische Grundbegriffe}
\label{b47guide:ch:foundations}

Dieser Leitfaden ergänzt Band~47 um eine geordnete Auswahl zentraler
Definitionen, Charakterisierungen und Folgerungen für metrische Räume.
Er enthält keine Beweise oder Beweisskizzen. Die Nummern mit Präfix $47H$
gehören zu dieser Ergänzung. Die Auswahl umfasst den metrischen Kern des
Hauptbands und weitere klassische Hauptresultate; sie beansprucht keine
vollständige Aufzählung aller bekannten oder ableitbaren Sätze.

\section{Voraussetzungen und Bandgrenzen}

Alle Räume sind metrische Räume, sofern nichts anderes angegeben ist.
Teilräume tragen stets die eingeschränkte Metrik. Zusätzliche Eigenschaften
wie Vollständigkeit oder Kompaktheit werden ausdrücklich genannt. Leere
Räume sind zugelassen; Aussagen, die einen Punkt verlangen, setzen
Nichtleerheit voraus. Der mengentheoretische Rahmen ist ZFC.

Behandelt werden Aussagen über Abstände und die durch sie bestimmten
Begriffe. Manche dieser Aussagen besitzen Verallgemeinerungen auf andere
Strukturen; hier werden ausschließlich ihre metrischen Fassungen verwendet.
Es werden keine topologischen, uniformen, normierten oder algebraischen
Räume als eigene axiomatische Theorie eingeführt.

\begin{center}
\small
\begin{tabularx}{\linewidth}{@{}p{.72\linewidth}X@{}}
\toprule
Eigenständiges Thema außerhalb dieses Leitfadens & Band\\
\midrule
Allgemeine Funktionen; Injektivität, Surjektivität, Bijektivität & 5--8\\
Allgemeine Äquivalenzrelationen & 11\\
Allgemeine Halbordnungen & 12\\
Allgemeine Schranken-, Infimums- und Supremumssätze & 13--14\\
Totale Ordnungen und Wohlordnungen & 15--16\\
Reelle Ordnungssätze und spezielle reelle Folgenresultate & 19, 21\\
\bottomrule
\end{tabularx}
\end{center}

Die reellen Zahlen dienen als Wertebereich der Metrik. Infimum und Supremum
werden nur zur Definition metrischer Größen und in metrisch bestimmten
Fortsetzungsformeln benutzt; ihre allgemeine
Theorie wird vorausgesetzt. Funktionen erscheinen mit metrischen
Voraussetzungen, etwa Stetigkeit, Abstandserhaltung oder Kontraktion.

\section{Metrik, Teilraum und Kugeln}

\begin{definition}[Metrik und metrischer Raum]
\label{b47guide:def:metric}
Eine \emph{Metrik} auf $X$ ist eine Funktion $d:X\times X\to\mathbb R$ mit
\[
\begin{aligned}
 d(x,y)&\geq0,&d(x,y)=0&\iff x=y,\\
 d(x,y)&=d(y,x),&d(x,z)&\leq d(x,y)+d(y,z)
 \qquad(x,y,z\in X).
\end{aligned}
\]
Das Paar $(X,d)$ heißt \emph{metrischer Raum}. Für $A\subseteq X$ heißt
$d_A=d|_{A\times A}$ die \emph{Teilraummetrik}.
\end{definition}

\begin{theorem}[Einschränkung und grundlegende Abstandsschranken]
Jeder Teilraum $(A,d_A)$ ist metrisch. Für $x,y,z,w\in X$ gelten
\[
 |d(x,z)-d(y,z)|\leq d(x,y),\qquad
 |d(x,y)-d(z,w)|\leq d(x,z)+d(y,w).
\]
Für jede endliche Punktfolge $x_0,\ldots,x_m$ gilt
$d(x_0,x_m)\leq\sum_{j=1}^{m}d(x_{j-1},x_j)$.
\end{theorem}

\begin{definition}[Offene und abgeschlossene Kugel]
Für $x\in X$ und $r>0$ seien
\[
 B_r(x)=\{y\in X:d(x,y)<r\},\qquad
 C_r(x)=\{y\in X:d(x,y)\leq r\}.
\]
Der Buchstabe $C$ bezeichnet die abgeschlossene Kugel; $\overline A$
bezeichnet im Folgenden den Abschluss einer Menge.
\end{definition}

\begin{theorem}[Kugelverkleinerung und Punkttrennung]
Ist $y\in B_r(x)$, so gilt
$B_{r-d(x,y)}(y)\subseteq B_r(x)$.
Sind $x\ne y$ und $0<r\leq d(x,y)/2$, so sind $B_r(x)$ und $B_r(y)$
disjunkt. Für $a\in A\subseteq X$ ist die Teilraumkugel
$B_r^{A}(a)=A\cap B_r^{X}(a)$.
\end{theorem}

\begin{definition}[Offen, abgeschlossen, Abschluss, dicht]
\label{b47guide:def:closure}
$U\subseteq X$ heißt \emph{offen}, wenn jeder Punkt $x\in U$ eine Kugel
$B_r(x)\subseteq U$ mit $r>0$ besitzt. $A\subseteq X$ heißt
\emph{abgeschlossen}, wenn $X\setminus A$ offen ist. Es seien
\[
 \overline A=\{x\in X:\forall r>0\ (B_r(x)\cap A\ne\varnothing)\},
 \qquad
 A^\circ=\{x\in A:\exists r>0\ (B_r(x)\subseteq A)\}.
\]
$\overline A$ heißt \emph{Abschluss}, $A^\circ$ \emph{Inneres}; $A$ heißt
\emph{dicht in $X$}, wenn $\overline A=X$.
\end{definition}

\begin{theorem}[Abgeschlossenheit von Kugeln und metrischer Abschluss]
Offene Kugeln sind offen, abgeschlossene Kugeln sind abgeschlossen.
$\overline A$ ist die kleinste abgeschlossene Teilmenge von $X$, die $A$
enthält. Es gilt
\[
 B_r(x)\subseteq\overline{B_r(x)}\subseteq C_r(x).
\]
Die zweite Inklusion kann echt sein. In einem Teilraum $Y\subseteq X$ gilt
für $A\subseteq Y$ die Formel $\overline A^{\,Y}=Y\cap\overline A^{\,X}$.
\end{theorem}

\section{Beschränktheit und Abstand zu einer Menge}

\begin{definition}[Beschränktheit und Durchmesser]
\label{b47guide:def:diameter}
$A\subseteq X$ heißt \emph{beschränkt}, wenn ein $M\geq0$ existiert mit
$d(a,b)\leq M$ für alle $a,b\in A$. Für nichtleeres $A$ sei
\[
 \operatorname{diam}A=\sup\{d(a,b):a,b\in A\}\in[0,\infty].
\]
Es gelte $\operatorname{diam}\varnothing=0$. Eine Folge heißt beschränkt,
wenn ihre Wertemenge beschränkt ist.
\end{definition}

\begin{theorem}[Kugel- und Abschlusskriterium der Beschränktheit]
Eine nichtleere Menge $A\subseteq X$ ist genau dann beschränkt, wenn sie in
einer Kugel endlichen Radius liegt. Bei Beschränktheit liegt sie für jeden
vorgegebenen Mittelpunkt $x\in X$ in einer genügend großen Kugel um $x$.
Es gilt $\operatorname{diam}\overline A=\operatorname{diam}A$.
\end{theorem}

\begin{definition}[Abstand von einem Punkt zu einer Menge]
Für $\varnothing\ne A\subseteq X$ und $x\in X$ sei
\[
 d(x,A)=\inf\{d(x,a):a\in A\}.
\]
Ein Punkt $a\in A$ mit $d(x,a)=d(x,A)$ heißt \emph{nächster Punkt} von $A$
zu $x$, sofern er existiert.
\end{definition}

\begin{theorem}[Abstandsfunktion und Abschlusserkennung]
\label{b47guide:thm:distance-set}
Für $\varnothing\ne A\subseteq X$ und $x,y\in X$ gelten
\[
 |d(x,A)-d(y,A)|\leq d(x,y),\qquad
 d(x,A)=d(x,\overline A),\qquad
 x\in\overline A\iff d(x,A)=0.
\]
Insbesondere erkennt $d(x,A)=0$ genau die Punkte von $A$, wenn $A$
abgeschlossen ist. Abgeschlossenheit allein sichert keinen nächsten Punkt.
\end{theorem}

\section{Metrische Konstruktionen und Vergleich}

Die folgenden Aussagen verwenden bereits die Begriffe Cauchy-Folge
(Definition~\ref{b47guide:def:cauchy}), Vollständigkeit
(Definition~\ref{b47guide:def:complete}), totale Beschränktheit
(Definition~\ref{b47guide:def:totally-bounded}) und Kompaktheit
(Definition~\ref{b47guide:def:compact}).

\begin{definition}[Gleichwertigkeit von Metriken]
Zwei Metriken $d,\rho$ auf derselben Menge heißen \emph{topologisch
äquivalent}, wenn sie dieselben offenen Mengen bestimmen. Sie heißen
\emph{gleichmäßig äquivalent}, wenn zu jedem $\varepsilon>0$ ein $\delta>0$
existiert, so dass für alle $x,y$ aus $d(x,y)<\delta$ die Ungleichung
$\rho(x,y)<\varepsilon$ folgt, und diese Bedingung auch mit vertauschten
Metriken gilt.
\end{definition}

\begin{theorem}[Beschränkte Ersatzmetrik]
\label{b47guide:thm:bounded-metric}
Zu jeder Metrik $d$ ist $d_b(x,y)=\min\{1,d(x,y)\}$ eine beschränkte,
gleichmäßig äquivalente Metrik. Sie hat dieselben konvergenten Folgen,
dieselben Cauchy-Folgen und dieselben total beschränkten Teilmengen wie $d$.
$(X,d_b)$ ist genau dann vollständig, wenn $(X,d)$ vollständig ist.
\end{theorem}

\begin{theorem}[Was ein Metrikwechsel erhält]
Topologisch äquivalente Metriken haben dieselben konvergenten Folgen mit
denselben Grenzwerten und dieselben kompakten Teilmengen. Umgekehrt
charakterisiert die Übereinstimmung aller Folgengrenzwerte ihre topologische
Äquivalenz. Gleichmäßige Äquivalenz erhält zusätzlich Cauchy-Folgen,
Vollständigkeit und totale Beschränktheit. Topologische Äquivalenz allein
erhält Vollständigkeit im Allgemeinen nicht.
\end{theorem}

